%% =============================================================================
%% WWW 2027 (ACM Web Conference) research track submission -- DRAFT SCAFFOLD
%% Template per CFP: \documentclass[sigconf, anonymous, review]{acmart}
%% Long paper: 8 content pages + references + optional appendix (12 pages max).
%% Track: "User Modeling, Personalization and Agentic Web Users" (chosen at submission).
%%
%% Build:  ~/.local/bin/tectonic main.tex     (from the paper/ directory)
%%
%% IMPORTANT (WWW'27 policy): LLMs may only help rephrase text; authors are
%% solely responsible for the content. This file is a structural scaffold written
%% with an assistant -- rewrite the prose in your own words before submission.
%% Red [..] = result placeholder. Blue [Note: ..] = guidance for the authors;
%% set \shownotesfalse to hide notes.
%% =============================================================================
\documentclass[sigconf,anonymous,review]{acmart}

\usepackage{booktabs}
\usepackage{multirow}
\usepackage{amsmath}
\usepackage{xcolor}

%% ---- Section heading case ----------------------------------------------------
%% acmart sigconf sets section titles in ALL CAPS. \plainheadingstrue keeps the
%% author's capitalisation. This modifies the official template: set it to
%% \plainheadingsfalse before submission unless the chairs confirm it's allowed.
\newif\ifplainheadings \plainheadingstrue
\ifplainheadings
  \makeatletter
  \def\@secfont{\sffamily\bfseries\section@raggedright}
  \makeatother
\fi

%% ---- Drafting macros ---------------------------------------------------------
\newif\ifshownotes \shownotestrue
\newcommand{\tbd}[1]{\textcolor{red}{[#1]}}
\newcommand{\note}[1]{\ifshownotes\textcolor{blue}{[Note: #1]}\fi}
\newcommand{\method}{\textsc{Hug}}      % working name -- rename freely
\newcommand{\figplaceholder}[2]{%
  \fbox{\parbox[c][#1][c]{0.95\linewidth}{\centering\small\textcolor{red}{#2}}}}
\newcommand{\R}{\mathbb{R}}
\newcommand{\bm}[1]{\boldsymbol{#1}}

%% ---- ACM metadata (placeholders; filled at camera-ready) --------------------
\setcopyright{acmlicensed}
\copyrightyear{2027}
\acmYear{2027}
\acmDOI{XXXXXXX.XXXXXXX}
\acmConference[WWW '27]{Proceedings of the ACM Web Conference 2027}{May 10--14, 2027}{Seoul, Republic of Korea}
\acmISBN{978-1-4503-XXXX-X/2027/05}

\begin{document}

\title{Shared or Specific? Evidence-Adaptive Sparse Fusion of Graph and Sequence Views for Click-Through Rate Prediction}
\note{Title is a working title. Alternatives: ``Separate, Then Fuse: \dots''.}

\author{Anonymous Author(s)}
\affiliation{\institution{Anonymous Institution}\country{}}

\begin{abstract}
Modern click-through rate (CTR) models draw on two complementary sources of
evidence: the relational structure linking users, items, creators and
categories, and the user's recent behaviour sequence. Graph encoders capture
the former and sequence transformers the latter, but the two views are almost
always fused by concatenation followed by an MLP. We argue that this is the
weak link. The views overlap substantially, so concatenation counts shared
evidence twice. It also lets the stronger sequence view dominate exactly
where the structural view is most valuable: long-tail items with little or no
interaction history. We propose \method{}, which keeps the two views separate
through encoding and then fuses them by decomposing each into sparse codes
over a \emph{shared} dictionary and \emph{view-private} dictionaries. The
number of active codes adapts to how much evidence each prediction rests on,
so sparse-evidence cases fall back to shared, prior-like codes. Under a
strictly causal evaluation protocol on three datasets (short video, news and
Q\&A), \method{} improves AUC by \tbd{x.xx\%} over the strongest tuned
baseline. The gains concentrate on cold items (\tbd{+x.xx} AUC in the
lowest-evidence bucket). The learned decomposition also quantifies how much
information the two views share (\tbd{xx\%}). We further show that two common
shortcuts, frozen graph encoders and item statistics aggregated over the whole
window, change conclusions: \tbd{frozen graph encoders underperform a
features-only control by x.xx AUC}.
\end{abstract}

\begin{CCSXML}
<ccs2012>
<concept>
<concept_id>10002951.10003317.10003347.10003350</concept_id>
<concept_desc>Information systems~Recommender systems</concept_desc>
<concept_significance>500</concept_significance>
</concept>
</ccs2012>
\end{CCSXML}
\ccsdesc[500]{Information systems~Recommender systems}

\keywords{CTR prediction, graph neural networks, sequential recommendation, multi-view fusion, sparse representations}

\maketitle

%% =============================================================================
\section{Introduction}
%% =============================================================================

Predicting whether a user will click an item is the core of web-scale
recommendation. Two kinds of evidence dominate modern CTR models. The first is
\emph{relational}. Users interact with items; items are made by creators and
tagged with categories. This structure is slow-changing, and graph encoders in
the LightGCN family exploit it well~\cite{he2020lightgcn,yu2022simgcl}. The
second is \emph{sequential}. The user's most recent actions reveal their
current intent, which transformer-based history encoders capture with target
attention~\cite{zhou2018din,chen2019bst,xia2023transact,zhai2024hstu}.

Both views are useful, and many systems combine them~\cite{chang2021surge,wang2020gcegnn,zhang2022gcl4sr}.
The way they are combined has received far less scrutiny than the encoders.
The default is to concatenate the view embeddings and let an MLP sort it out.
We argue that this default is the weak link, for two reasons.

\emph{The views overlap.} An item's graph neighbourhood and the sequences it
appears in are generated by the same users clicking the same items, so much
of what one view encodes the other encodes too. An MLP over the concatenation
has no inductive pressure to separate redundant evidence from complementary
evidence, and it double-counts the former.

\emph{The views are unequally reliable, and the imbalance shifts with the
amount of evidence.} When a candidate item has a rich history, the sequence
view is strong and the structural view adds little. For long-tail items,
which make up \tbd{58\% of validation impressions on KuaiRand-1K}
\note{from the Spec 01 cold bucket; re-measure on the final protocol}, the
sequence view has little to say about the candidate. The structural view is
then the main route by which signal reaches the item: through its creator,
its categories and co-engaging users. A fusion that weights views the same
way for every input cannot express this shift.

These observations suggest a different design principle: \emph{keep the
views separate long enough to measure what each contributes, then fuse them
in a way that is explicit about shared versus view-specific information and
that adapts to the available evidence.} \method{} implements this principle.
Each view is encoded by a standard, strong encoder: a relation-aware LightGCN
with contrastive regularisation for structure, and a target-attention
transformer for sequence. The novelty lies in the fusion. Each view's
representation is projected onto sparse codes over a dictionary shared by both
views plus a private dictionary per view. A cross-view alignment objective
makes the shared codes agree; a decorrelation objective keeps the private codes
complementary. The number of active codes is governed by an evidence-dependent
budget: predictions resting on little evidence are pushed towards a few shared
atoms, while well-supported predictions may use many private atoms. The
downstream MLP thus receives a representation that is sparse, separates
redundant from complementary evidence, and is calibrated to how much the model
actually knows.

We deliberately use standard encoders. This isolates the contribution: every
fusion variant we compare sees identical encoder outputs, so differences are
attributable to fusion alone. The design also yields a quantity that
concatenation cannot provide: the fraction of predictive information the two
views share. That fraction can be reported, and it can be analysed by evidence
level.

A second, methodological thread runs through the paper. Graph-based CTR
pipelines are unusually easy to leak: interaction edges, item popularity
statistics and session structure computed over the full log all contain
future information~\cite{ji2023leakage}. We evaluate under a strict
point-in-time protocol, in which every input for an impression at time $t$
uses only interactions before $t$. We pair it with a features-only control.
We show that this matters: \tbd{under the causal protocol, a frozen
(randomly initialised) graph encoder underperforms the features-only control
by x.xx AUC, and whole-window item statistics inflate a sequential baseline's
training-time AUC while collapsing its validation AUC}.

\paragraph{Contributions.}
\begin{itemize}
  \item We identify concatenation-based fusion of graph and sequence views as a
        bottleneck in CTR prediction. We show empirically that it underuses the
        structural view on long-tail items \tbd{evidence: Fig.~\ref{fig:motivation}}.
  \item We propose \method{}, an evidence-conditioned sparse fusion layer.
        It encodes the graph and sequence views as sparse codes over a shared
        dictionary and view-private dictionaries. Each view's budget of private
        atoms is set per impression from the interaction evidence available
        \emph{to that view}, so poorly supported views are forced onto atoms
        both views agree on (\S\ref{sec:fusion}). To our knowledge, this is the
        first fusion layer for CTR whose shared/private allocation is tied to
        per-impression evidence. Because the codes are sparse, the same mechanism
        yields a countable, per-prediction measure of view overlap.
  \item We describe and release a causal evaluation protocol for graph-based CTR
        models: time-masked graph snapshots, point-in-time features, a
        features-only control, and enforced test-set isolation. We show how common
        shortcuts change conclusions (\S\ref{sec:protocol}, RQ5).
  \item On three datasets (short video, news, Q\&A), \method{} outperforms tuned sequential,
        feature-interaction and graph-based CTR baselines
        \tbd{by x.xx--x.xx AUC}, with gains concentrated where the mechanism predicts.
        The learned decomposition reveals \tbd{finding about shared vs.\ private
        information across evidence levels}.
\end{itemize}

\begin{figure}[t]
  \centering
  \figplaceholder{3.6cm}{Motivation figure: AUC by candidate-item evidence bucket
  (0 / 1--4 / $\geq$5 prior interactions) for sequence-only, graph-only,
  concat fusion, and \method{}. Expected shape: concat $\approx$ max(single views)
  in the cold bucket, while \method{} exceeds both.}
  \caption{\tbd{Concatenation fusion fails to combine the views where evidence
  is scarce.} \note{Populate from heavy-run arms N1, N2, N4 plus the \method{} run.}}
  \label{fig:motivation}
\end{figure}

%% =============================================================================
\section{Related Work}
%% =============================================================================

\paragraph{Graph-based recommendation.}
Graph neural networks propagate collaborative signal over user--item graphs.
LightGCN~\cite{he2020lightgcn} showed that removing feature transforms and
nonlinearities improves collaborative filtering, and contrastive variants
(SimGCL, XSimGCL~\cite{yu2022simgcl,yu2023xsimgcl}) further help long-tail
items. Heterogeneous and knowledge-aware models
(R-GCN, KGAT, HGT~\cite{schlichtkrull2018rgcn,wang2019kgat,hu2020hgt})
model typed relations. Benchmarking studies, however, find that well-tuned
simple models are hard to beat~\cite{lv2021simplehgn,dacrema2019progress,rendle2020ncf}.
Precomputed, parameter-free propagation (SGC, SIGN~\cite{wu2019sgc,frasca2020sign})
can rival trained message passing. This motivates our frozen-encoder control.

\paragraph{Sequential CTR models.}
DIN~\cite{zhou2018din} introduced target attention over behaviour histories.
BST, SASRec and TransAct~\cite{chen2019bst,kang2018sasrec,xia2023transact}
use transformers, and HSTU~\cite{zhai2024hstu} and
WuKong~\cite{zhang2024wukong} study scaling. Feature-interaction models such
as DCNv2 and Fi-GNN~\cite{wang2021dcnv2,li2019fignn} model cross features
without user--item graphs.

\paragraph{Combining graphs and sequences.}
SR-GNN and GCE-GNN~\cite{wu2019srgnn,wang2020gcegnn} build graphs from
sessions. SURGE~\cite{chang2021surge} constructs interest graphs from
sequences. GCL4SR, MCLSR and MVCrec~\cite{zhang2022gcl4sr,wang2022mclsr,zhou2026mvcrec}
align graph and sequence views with contrastive objectives, and
SelfGNN~\cite{liu2024selfgnn} combines graph and sequential self-supervision.
DCRec~\cite{yang2023dcrec} unifies sequential patterns with a global
collaborative graph through cross-view contrast, using popularity-aware
augmentation weights. In CTR, graph-enhanced interest networks target users
with sparse behaviour sequences~\cite{liu2025geinsparse}. These works differ
mainly in how each view is \emph{encoded} and aligned. Fusion is usually
concatenation, attention or a learned gate. Our focus is the fusion layer
itself: what is shared, what is view-specific, and how the allocation between
the two should depend on each view's evidence.

\paragraph{Shared/private decomposition.}
Separating view-invariant from view-specific factors is established in
multimodal learning (MISA~\cite{hazarika2020misa}; the multi-view information
bottleneck~\cite{federici2020mib}) and has been carried into recommendation.
CMDL~\cite{lin2025cmdl} and MRdIB~\cite{wang2025mrdib} disentangle modalities,
DisenCDR~\cite{cao2022disencdr} and CDRIB~\cite{cao2022cdrib} separate
domain-shared from domain-specific user factors, and RIDRec~\cite{xia2026ridrec}
decomposes the predictive roles of semantic and collaborative signals into
shared and unique parts, arguing against concatenation. Shared and private
\emph{sparse dictionaries} date back to factorised latent spaces with
structured sparsity~\cite{jia2010factorized}. Recent group-sparse autoencoders
study when multimodal dictionaries truly share atoms~\cite{kaushik2026groupsae}.
\method{} adopts this decomposition with sparse codes, applies it to the
graph and sequence views of a CTR model, and, unlike prior work, makes the
shared/private allocation depend on evidence.

\paragraph{Evidence-adaptive capacity.}
Tying representation capacity to popularity is a mature idea in
recommendation. Mixed-dimension embeddings~\cite{ginart2021mde}, AutoEmb~\cite{zhao2021autoemb}
and PEP~\cite{liu2021pep} size embeddings by frequency or learned pruning.
VBAE~\cite{zhu2022vbae} fuses a collaborative and a feature view through a
user-dependent information channel whose bandwidth reflects how sufficient
the collaborative evidence is. GateSID~\cite{zhu2026gatesid} gates semantic and
collaborative signals by item maturity, with stronger alignment for cold items.
\method{} differs in where and how capacity adapts. Adaptation happens per
impression inside the fusion layer, acts \emph{asymmetrically} on private
versus shared atoms, and is driven by a separate evidence signal for each view.
Throughout, ``evidence'' means an interaction-count proxy
(\S\ref{sec:fusion}), not the Dirichlet evidence of evidential deep learning.
Sparse coding tools~\cite{makhzani2014ksparse,gao2024sae,louizos2018l0} and
information-bottleneck objectives~\cite{tishby2000ib,alemi2017vib} underpin the
mechanism. In recommendation, discrete codes appear mainly as item
tokenisers~\cite{rajput2023tiger}, and embedding collapse has been studied at
scale~\cite{guo2024collapse}.

\paragraph{Evaluation rigour.}
Data leakage in offline recommender evaluation is well documented~\cite{ji2023leakage},
as are tuning disparities between proposed models and
baselines~\cite{dacrema2019progress,rendle2020ncf,zhu2021bars}. CTR models
with ID embeddings are known to overfit after one epoch~\cite{zhang2022oneepoch}.
Our protocol extends point-in-time evaluation to graph inputs.

%% =============================================================================
\section{Problem Setting and Causal Protocol}
\label{sec:protocol}
%% =============================================================================

\paragraph{Data.}
Let $\mathcal{U}, \mathcal{V}, \mathcal{A}, \mathcal{C}$ denote users, items,
creators (authors) and categories. The log is a time-ordered set of impressions
$\mathcal{L} = \{(u, v, t, y, \bm{x})\}$, where $y \in \{0,1\}$ indicates a
click and $\bm{x}$ holds context (e.g., page tab, hour). The task is to estimate
$P(y = 1 \mid u, v, t, \mathcal{H}_{<t})$, where $\mathcal{H}_{<t}$ is the
information available strictly before $t$.

\paragraph{Point-in-time rule.}
Every input for an impression at time $t$ is computed from interactions with
timestamp $< t$. This covers graph edges, item statistics, behaviour histories
and ID vocabularies. Model weights are fit on a training period that strictly
precedes validation and test, using a global chronological split
(\tbd{70/10/20}). Ties never straddle a boundary.

\paragraph{Time-masked graph snapshots.}
We build one heterogeneous graph with relations $\mathcal{R}$: user--item
feedback (click, like, dislike, comment, forward), item--creator,
item--category and item--item transitions, plus reverse relations. Every
behavioural edge carries its timestamp. The timeline is cut at boundaries
$b_0 < b_1 < \dots$ (daily, with the split points inserted). An impression in
$[b_k, b_{k+1})$ is encoded from snapshot $\mathcal{G}_k$, which contains only
edges with time $< b_k$; degree normalisation is computed from $\mathcal{G}_k$
alone. Item popularity statistics are computed per row from strictly earlier
interactions.

\paragraph{Controls and test isolation.}
Every comparison includes a \emph{features-only} control: the same head and
inputs with no graph or sequence encoder. Hyperparameters are selected on
validation with an equal trial budget for every model. Test is scored once
per reported configuration, and the harness enforces and logs this.

%% =============================================================================
\section{Method}
%% =============================================================================

\method{} (\emph{Heterogeneous Unified Graph}) is built on a single
heterogeneous, time-stamped graph: users, items, creators and categories,
linked by every feedback type and by in-session transitions. This graph is
\emph{unified} in the sense that it is the one source of truth for both
views. The graph view propagates over its time-masked snapshots, and the
sequence view reads behaviour histories from the same timestamped
interactions under the same point-in-time rule. The model is deliberately
not unified: the views stay separate until fusion.
Figure~\ref{fig:arch} gives an overview. \method{} has three stages:
(i) a shared input layer, (ii) two view encoders run in parallel, and
(iii) evidence-adaptive sparse fusion followed by a prediction head.

\begin{figure*}[t]
  \centering
  \figplaceholder{4.2cm}{Architecture figure: input layer $\rightarrow$
  [graph view: relation-aware LightGCN on snapshot $\mathcal{G}_k$] and
  [sequence view: history transformer + target attention] $\rightarrow$
  sparse encoders onto shared dictionary $D^{\mathrm{sh}}$ and private
  dictionaries $D^{\mathrm{G}}, D^{\mathrm{S}}$ with evidence-dependent budget
  $\rightarrow$ MLP head.}
  \caption{Overview of \method{}.}
  \label{fig:arch}
\end{figure*}

\subsection{Input layer}
Each node receives an initial embedding
$\bm{h}^{(0)} \in \R^d$ that combines an ID embedding with projected features.
ID vocabularies are built from the training period only; IDs seen fewer than
$m$ times share an out-of-vocabulary row. Creators and categories are
represented by ID embeddings.

\subsection{Graph view}
We use a relation-aware variant of LightGCN. For layer $l$,
\begin{equation}
  \bm{h}^{(l+1)}_i = \sum_{r \in \mathcal{R}} \pi^{(l)}_r
    \sum_{j \in \mathcal{N}_r(i)} \frac{\bm{h}^{(l)}_j}{\sqrt{|\mathcal{N}_r(i)|\,|\mathcal{N}_r(j)|}},
\end{equation}
with $\bm{\pi}^{(l)} = \mathrm{softmax}(\bm{a}^{(l)})$,
where neighbourhoods come from snapshot $\mathcal{G}_k$ and $\bm{a}^{(l)}$
are learned relation logits. The graph representation is the layer mean
$\bm{g}_i = \frac{1}{L+1}\sum_{l=0}^{L} \bm{h}^{(l)}_i$. Following
XSimGCL~\cite{yu2023xsimgcl}, we add a cross-layer contrastive term
$\mathcal{L}_{\mathrm{cl}}$ over the users and items in each batch. The
graph-view input to fusion is
$\bm{e}^{\mathrm{G}} = [\bm{g}_u \,\|\, \bm{g}_v \,\|\, \bm{g}_u \odot \bm{g}_v]$.

\subsection{Sequence view}
For an impression $(u, v, t)$, the history $S_u(t)$ is the user's last $N$
clicked items before $t$. Each token is the item's input embedding plus
embeddings of recency rank, log-bucketed time gap and a same-session flag. A
light transformer encodes the tokens. The candidate item queries the outputs
by target attention, giving $\bm{e}^{\mathrm{S}}$.

\subsection{Evidence-adaptive sparse fusion}
\label{sec:fusion}
\note{The method is specified (docs/specs/05-sparse-fusion.md) but not yet
implemented. Finalise the equations once validation selects between the top-$k$
schedule and hard-concrete gates.}

\paragraph{Per-view evidence.}
Each view's reliability depends on the support it has for this particular
prediction. We define, from interactions strictly before $t$,
\begin{align}
  n^{\mathrm{G}} &= \log(1+c_v) + \log(1+c_{a(v)}), \\
  n^{\mathrm{S}} &= \log(1+|S_u(t)|) + \log(1+c_v),
\end{align}
where $c_v$ and $c_{a(v)}$ are the prior interaction counts of the candidate
item and of its creator, and $|S_u(t)|$ is the user's history length. A new
video by an active creator thus has strong graph support but weak sequence
support. This asymmetry is what lets the structural view carry cold items.

\paragraph{Shared and private dictionaries.}
Each view's input is projected to $\bm{p}^{\nu} = \mathrm{LN}(W^{\nu}\bm{e}^{\nu}) \in \R^{d'}$,
$\nu \in \{\mathrm{G}, \mathrm{S}\}$. We learn a shared dictionary
$D^{\mathrm{sh}} \in \R^{m_s \times d'}$ and private dictionaries
$D^{\mathrm{G}}, D^{\mathrm{S}} \in \R^{m_p \times d'}$, with unit-norm atoms.
An encoder maps each view to non-negative codes over the shared atoms and its
own private atoms,
$[\bm{a}^{\nu,\mathrm{sh}} \,\|\, \bm{a}^{\nu,\mathrm{pr}}] = \mathrm{ReLU}(E^{\nu}\bm{p}^{\nu} + \bm{b}^{\nu})$.

\paragraph{Evidence-dependent sparsity.}
We keep the top $k_{\mathrm{sh}}$ shared codes for every impression. The
number of private codes kept for view $\nu$ grows with that view's evidence:
\begin{equation}
  k^{\nu}_{\mathrm{pr}} = k_{\min} + \Big\lfloor (k_{\max} - k_{\min})\,
     \min\!\big(1, n^{\nu}/n^{\nu}_{\mathrm{ref}}\big) \Big\rceil,
\end{equation}
where $n^{\nu}_{\mathrm{ref}}$ is the 95th percentile of $n^{\nu}$ over
training impressions. Poorly supported views are therefore restricted to a few
private atoms and must express themselves through atoms shared with the other
view. Read through the information bottleneck~\cite{tishby2000ib,alemi2017vib},
the schedule is a per-impression, per-view rate constraint. A soft alternative
places hard-concrete gates~\cite{louizos2018l0} on the codes, with an expected-$L_0$
penalty on private atoms weighted by $\lambda^{\nu}(n^{\nu}) = \lambda_0/(1+n^{\nu})$.
\tbd{Report whichever validates better; the other goes in the appendix.}

\paragraph{Objectives.}
Shared codes from the two views should describe the same evidence, and private
codes should carry what the other view lacks:
\begin{align}
  \mathcal{L}_{\mathrm{align}} &= \mathrm{InfoNCE}\big(\bm{a}^{\mathrm{G},\mathrm{sh}}, \bm{a}^{\mathrm{S},\mathrm{sh}}\big), \\
  \mathcal{L}_{\mathrm{dec}}   &= \big\lVert \mathrm{Cov}\big(\bm{a}^{\mathrm{G},\mathrm{pr}}, \bm{a}^{\mathrm{S},\mathrm{pr}}\big) \big\rVert_F^2 / m_p^2, \\
  \mathcal{L}_{\mathrm{rec}}   &= \textstyle\sum_{\nu} \big\lVert \mathrm{sg}(\bm{p}^{\nu}) - D^{\mathrm{sh}\top}\bm{a}^{\nu,\mathrm{sh}} - D^{\nu\top}\bm{a}^{\nu,\mathrm{pr}} \big\rVert_2^2 .
\end{align}
InfoNCE~\cite{oord2018infonce} uses in-batch negatives. $\mathcal{L}_{\mathrm{dec}}$
is the linear-kernel case of HSIC~\cite{gretton2005hsic}. The reconstruction
target is detached ($\mathrm{sg}$), so the projections cannot collapse towards
whatever is easiest to reconstruct.

\paragraph{Fused representation and head.}
The head receives
$\bm{z} = [\tfrac{1}{2}(\bm{a}^{\mathrm{G},\mathrm{sh}} + \bm{a}^{\mathrm{S},\mathrm{sh}})
 \,\|\, \bm{a}^{\mathrm{G},\mathrm{pr}} \,\|\, \bm{a}^{\mathrm{S},\mathrm{pr}} \,\|\, \bm{x}]$
and outputs a logit through an MLP. The training objective is
\begin{equation}
  \mathcal{L} = \mathcal{L}_{\mathrm{BCE}} + \alpha\,\mathcal{L}_{\mathrm{cl}}
   + \beta\,\mathcal{L}_{\mathrm{align}} + \gamma\,\mathcal{L}_{\mathrm{dec}}
   + \delta\,\mathcal{L}_{\mathrm{rec}}.
\end{equation}
The shared-mass ratio
$\rho = \sum_{\nu}\lVert \bm{a}^{\nu,\mathrm{sh}}\rVert_1 / \sum_{\nu}\lVert \bm{a}^{\nu}\rVert_1$
gives a per-prediction measure of view overlap. We also report the share of
shared atoms that both views actually co-activate, which guards against shared
atoms degenerating into private ones~\cite{kaushik2026groupsae}.

\subsection{Complexity}
\tbd{Per-step cost of graph propagation (sparse matrix products over
$|\mathcal{E}_k|$ edges), the transformer ($O(N^2 d)$ per row) and fusion
($O((m_s + m_g + m_q) d')$ per row); measured wall-clock in Appendix.}

%% =============================================================================
\section{Experiments}
%% =============================================================================

We address five research questions.
\textbf{RQ1:} Does \method{} outperform strong, equally tuned baselines?
\textbf{RQ2:} Is the gain attributable to fusion?
\textbf{RQ3:} Where do the gains come from?
\textbf{RQ4:} What do the shared and private codes capture?
\textbf{RQ5:} How much do protocol shortcuts change conclusions?

\subsection{Setup}

\paragraph{Datasets.}
KuaiRand-1K~\cite{gao2022kuairand} is a short-video log from Kuaishou
containing 1{,}000 users and their impressions, with creator, category and
multi-type feedback. MIND~\cite{wu2020mind} is a news-recommendation log with
impressions, categories, subcategories and Wikidata entities. Articles churn
daily, so most test impressions involve items never seen in training.
ZhihuRec~\cite{hao2021zhihurec} is a Q\&A log with answers linked to authors,
questions and topics, in which most test impressions come from users unseen in
training. The three datasets cover three domains and both cold-start axes. For
MIND and ZhihuRec we sample users at random (\tbd{20\% and 12.5\%}) and keep each
sampled user's full log. ZhihuRec logs click times separately from impression
times, so every label-derived input uses only labels \emph{known} before $t$.
Table~\ref{tab:data} summarises the splits.

\begin{table}[t]
  \caption{Dataset statistics after preprocessing (global chronological split).}
  \label{tab:data}
  \small
  \setlength{\tabcolsep}{3pt}
  \begin{tabular}{lrrr}
    \toprule
    & KuaiRand-1K & MIND & ZhihuRec \\
    \midrule
    Users              & \tbd{1,000} & \tbd{} & \tbd{} \\
    Items              & \tbd{4.37M} & \tbd{} & \tbd{} \\
    Train / val / test (M) & 8.19 / 1.17 / 2.34 & \tbd{} & \tbd{} \\
    Click rate         & 0.379 & \tbd{} & \tbd{} \\
    Cold-item test share & \tbd{} & \tbd{} & \tbd{} \\
    Cold-user test share & \tbd{} & \tbd{} & \tbd{} \\
    \bottomrule
  \end{tabular}
  \note{Re-verify every figure from the final preprocessing logs.}
\end{table}

\paragraph{Baselines.}
Sequential CTR: TransAct~\cite{xia2023transact}. Scaling and feature
interaction: WuKong~\cite{zhang2024wukong}\tbd{, DCNv2~\cite{wang2021dcnv2}}.
GNN-based CTR: Fi-GNN~\cite{li2019fignn}. All baselines use the FuxiCTR
implementations~\cite{zhu2021bars} with a pinned model zoo. Every baseline is
given the same row inputs as \method{}: identical behaviour histories,
point-in-time item statistics, context and categorical features. Each model,
including \method{}, receives \tbd{16} tuning trials on validation.

\paragraph{Fusion variants (RQ2).}
All fusion variants use identical, jointly trained encoders and differ only in
fusion: concatenation + MLP; gated fusion; an \emph{evidence-conditioned} gate
(GateSID-style~\cite{zhu2026gatesid}, the gate sees $n^{\mathrm{G}}, n^{\mathrm{S}}$);
mixture-of-experts routing~\cite{ma2018mmoe}; dense shared/private
decomposition with the same losses (MISA/CMDL-style~\cite{hazarika2020misa,lin2025cmdl});
\method{} with a constant sparsity budget; and full \method{}. The
evidence-conditioned gate is the key competitor: it tests whether evidence
adaptivity alone, without sparse shared/private codes, explains the gains.

\paragraph{Metrics.}
AUC, AP, LogLoss and per-user nDCG@10. We report means and standard deviations
over \tbd{3} seeds. Significance uses a per-user paired bootstrap
(1{,}000 resamples).

\paragraph{Implementation.}
\tbd{Embedding size, layers, dictionary sizes, optimiser, batch size, early
stopping; hardware (2$\times$ 95\,GB GPUs).} Code and protocol are released
at \tbd{anonymous link}.

\subsection{RQ1: Overall performance}

\begin{table}[t]
  \caption{Test performance (mean $\pm$ std over seeds). Best in bold;
  $^{\dagger}$ significant vs.\ best baseline ($p<0.05$).}
  \label{tab:main}
  \small
  \setlength{\tabcolsep}{4pt}
  \begin{tabular}{lcccc}
    \toprule
    Model & AUC & AP & LogLoss & nDCG@10 \\
    \midrule
    Features-only control & \tbd{} & \tbd{} & \tbd{} & \tbd{} \\
    TransAct      & \tbd{} & \tbd{} & \tbd{} & \tbd{} \\
    WuKong        & \tbd{} & \tbd{} & \tbd{} & \tbd{} \\
    Fi-GNN        & \tbd{} & \tbd{} & \tbd{} & \tbd{} \\
    \tbd{DCNv2}   & \tbd{} & \tbd{} & \tbd{} & \tbd{} \\
    \midrule
    Sequence view only & \tbd{} & \tbd{} & \tbd{} & \tbd{} \\
    Graph view only    & \tbd{} & \tbd{} & \tbd{} & \tbd{} \\
    Both, concat fusion & \tbd{} & \tbd{} & \tbd{} & \tbd{} \\
    \textbf{\method{}} & \tbd{} & \tbd{} & \tbd{} & \tbd{} \\
    \bottomrule
  \end{tabular}
\end{table}

\tbd{Narrative: \method{} vs.\ best baseline; whether the graph view adds over
the sequence-only model; the features-only control as a floor.}

\subsection{RQ2: Is the gain due to fusion?}
\tbd{Table: fusion variants with identical encoders. Key comparisons:
concat vs.\ \method{} (total effect of fusion); dense vs.\ sparse
shared/private (effect of sparsity); fixed vs.\ adaptive budget (effect of
evidence adaptivity).}

\subsection{RQ3: Where do the gains come from?}
\begin{figure}[t]
  \centering
  \figplaceholder{3.4cm}{AUC by candidate-item evidence (0 / 1--4 / $\geq$5)
  and by user history length (0 / 1--9 / 10--49 / 50), for the best baseline,
  concat fusion, and \method{}.}
  \caption{Gains by evidence level.}
  \label{fig:buckets}
\end{figure}
\tbd{Expected: the largest gains appear in low-evidence buckets, matching the
budget mechanism.}

\subsection{RQ4: What do shared and private codes capture?}
\tbd{(a) Shared-mass ratio $\rho$ vs.\ evidence level: does low evidence shift
mass to shared atoms? (b) Overall view overlap: average $\rho$, compared with
CKA between view embeddings. (c) Atom interpretability: top creators and
categories per atom. (d) Learned relation weights $\bm{\pi}^{(l)}$.}

\subsection{RQ5: Do protocol shortcuts change conclusions?}
\tbd{Table: (i) frozen random graph encoder vs.\ trained vs.\ features-only;
(ii) whole-window vs.\ point-in-time item statistics for a sequential
baseline; (iii) effect on the ranking of methods. Smoke-scale evidence so far
suggests that frozen encoders fall below the features-only control; confirm
at full scale.}

%% =============================================================================
\section{Conclusion}
%% =============================================================================
We argued that fusing graph and sequence views is a first-class design
problem in CTR prediction, not an afterthought. \method{} keeps the views
separate through encoding, then fuses them through sparse shared and private
codes whose budget adapts to the available evidence. \tbd{Summary of gains and
the overlap finding.} Our causal protocol and controls show that
\tbd{protocol finding}. Limitations: \tbd{datasets with dense random exposure;
online evaluation; dictionary size sensitivity}.

%% =============================================================================
\section*{Ethical Considerations}
This work uses publicly released, anonymised interaction logs under their
licences. No new user data were collected. Improved long-tail recommendation
may affect creator exposure; we report results by item popularity to make such
effects visible.

\bibliographystyle{ACM-Reference-Format}
\bibliography{references}

\appendix
\section{Protocol Details}
\tbd{Snapshot construction, tie handling, ID vocabulary, holdout diagnostic,
test-access logging.}

\section{Hyperparameters and Tuning Budgets}
\tbd{Search spaces and selected configurations for every model.}

\section{Additional Results}
\tbd{Validation tables, per-seed results, run time and memory.}

\end{document}
